\section{FHE Applications}


For FHE, the NTT must handle larger polynomial degrees ($n = 1024, 2048, 4096$) and
larger moduli ($q$ up to 60 bits for RNS representation). The techniques are the same,
but:

\begin{itemize}[nosep]
  \item \textbf{Larger degree}: More butterfly layers ($\log_2 n$ layers). AVX2 benefit
    increases because more parallelism is available.
  \item \textbf{Larger modulus}: 16-bit Montgomery does not suffice. We use 32-bit
    Montgomery with 64-bit intermediate products. AVX2 \texttt{VPMULLD} handles 32-bit
    multiplies (8 per instruction instead of 16).
  \item \textbf{RNS}: The Residue Number System represents large integers as tuples of
    residues modulo small primes. Each NTT operates on one residue---enabling parallelism
    across RNS channels.
\end{itemize}

%% ============================================================
